\section{Restrições e extensões}
Nessa seção, a menos que explicitamente indicado, consideraremos que $(M, d)$ e $(N, d')$ são espaços métricos arbitrários.

Os limites são preservados por restrições sempre que ainda fazem sentido.

A demonstração abaixo se baseia no fato de que se $f:M\rightarrow N$ é uma função e $i: A\rightarrow M$ é a inclusão de $A$ em $M$, então $f|_A = f\circ i$.
De fato, $\dom(f|_A)=\dom(f\circ i)=A$ e, dado $x \in A$, temos $f|_A(x)=f(x)=f(i(x))=f\circ i(x)$.

\begin{proposition}Sejam
    \begin{itemize}
        \item $A\subseteq M$ um subespaço,
        \item $f: M \to N$ uma função,
        \item $p$ ponto de acumulação de $A$,
        \item $f: M \to N$ uma função, e
        \item $L\in N$.
    \end{itemize}

    Se $L$ é limite de $f$ em $p$, então $L$ é limite de $f|_A$ em $p$.

    Em outras palavras, se existe $\lim_{x\to p} f(x)$, então existe $\lim_{x\to p} f|_A(x)$, e
    \[
        \lim_{x\to p} f|_A(x) = \lim_{x\to p} f(x).
    \]
\end{proposition}

\begin{proof}
    Suponha que $L$ é limite de $f$ em $p$.
    Seja $i:A\rightarrow M$ a inclusão de $A$ em $M$.

    Temos que $i$ é contínua em $p$ e $p$ é ponto de acumulação de $A$, então $p=\lim_{x\to p} i(x)$.
    Além disso, $i:A\rightarrow M$ é uma função injetora.

    Logo, pelo teorema da mudança de variável em limites,
    \[
        \lim_{x\to p} f|_A(x) = \lim_{x\to p} f(i(x)) = \lim_{x\to p} f(x) = L.
    \]
    E, portanto, $L$ é limite de $f|_A$ em $p$.    
\end{proof}

Disso, decorre que a continuidade também é preservada por restrições.

\begin{corollary}Sejam
    \begin{itemize}
        \item $A\subseteq M$ um subespaço,
        \item $f: A \to N$ uma função,
        \item $p\in A$, e
        \item $f: A \to N$ uma função.
    \end{itemize}

    Se $f$ é contínua em $p$, então $f|_A$ é contínua em $p$.

    Consequentemente, se $f$ é contínua, então $f|_A$ é contínua.
\end{corollary}

\begin{proof}
    Suponha que $f$ é contínua em $p$.

    Se $p$ é um ponto isolado de $A$, então $f|_A$ é contínua em $p$ trivialmente.

    Se $p$ é um ponto de acumulação de $A$, então, pela proposição anterior,

    \[
        \lim_{x\to p} f|_A(x) = \lim_{x\to p} f(x) = f(p),
    \]

    Assim, $f|_A$ é contínua em $p$.
\end{proof}

Dessa forma, se $f$ é contínua em $p$, então $f|_A$ é contínua em $p$.

A recíproca não é verdadeira!
No Exemplo~\ref{ex:restricaoContinuidade}, vimos que existem funções $f$ que não são contínuas em $p$, mas cuja restrição $f|_A$ é contínua em $p$.

Por outro lado, em alguns casos, a continuidade é preservada por extensões.

\begin{proposition}Sejam
    \begin{itemize}
        \item $A\subseteq M$ um subconjunto
        \item $f: M \to N$ uma função,
        \item $p\in A$ é ponto de acumulação de $A$,
        \item $f: M \to N$ uma função, e
        \item $L\in N$.
    \end{itemize}

    Suponha que:
        \begin{itemize}
            \item $f|_A$ é contínua em $p$, e
            \item $L$ é limite de $f|_A$ em $p$.
        \end{itemize}

    Se $L$ é limite de $f|_A$ em $p$, então $L$ é limite de $f$ em $p$.

    Em outras palavras, se existe $\lim_{x\to p} (f|_A)(x)$, então existe $\lim_{x\to p} f(x)$, e
    \[
        \lim_{x\to p} f|_A(x) = \lim_{x\to p} f(x).
    \]
\end{proposition}

\begin{proof}
    Suponha que $L$ é limite de $f|_A$ em $p$.
    Devemos mostrar que $L$ é limite de $f$ em $p$.
    Ou seja, dado $\epsilon>0$, devemos mostrar que existe $\delta>0$ tal que, para todo $x\in M\setminus\{p\}$,
    
    \[
        \text{se } d(x, p) < \delta \qquad\text{então}\qquad d'(f(x), L) < \epsilon.
    \]

    Sabemos que $L$ é limite de $f|_A$ em $p$, então, para o $\epsilon>0$ dado, existe $\delta_1>0$ tal que, para todo $x\in A\setminus\{p\}$,
    \[
        \text{se } d(x, p) < \delta_1 \qquad\text{então}\qquad d'((f|_A)(x), L) < \epsilon.
    \]

    Como $A$ é aberto e $p$ é um ponto de acumulação de $A$, existe $\delta_2>0$ tal que $B_M(p, \delta_2) \subseteq A$.

    Seja $\delta = \min\{\delta_1, \delta_2\}$.
    Fixe $x \in M\setminus\{p\}$ tal que $d(x, p) < \delta$.
    Assim, $d(x, p) < \delta_1$ e $d(x, p) < \delta_2$.
    
    Como $d(x, p) < \delta_2$, temos $x \in B_M(p, \delta_2) \subseteq A$.
    Assim, $x \in A\setminus\{p\}$ e $d(x, p) < \delta_1$, então $d'((f|_A)(x), L) < \epsilon$.
    Como $(f|_A)(x) = f(x)$, segue que $d'(f(x), L) < \epsilon$.
\end{proof}

\begin{corollary}Sejam
    \begin{itemize}
        \item $A\subseteq M$ um subespaço aberto,
        \item $f: A \to N$ uma função,
        \item $p\in A$, e
        \item $f: A \to N$ uma função.
    \end{itemize}

    Se $f|_A$ é contínua em $p$, então $f$ é contínua em $p$.

    Consequentemente, se $f|_A$ é contínua, então $f$ é contínua.
\end{corollary}
\begin{proof}
    Suponha que $f|_A$ é contínua em $p$.

    Se $p$ é um ponto isolado de $M$, então $f$ é contínua em $p$ trivialmente.

    Se $p$ é um ponto de acumulação de $M$, então é um ponto de acumulação de $A$, pois se $\{p\}$ fosse aberto em $A$, também seria em $M$, uma vez que $A\subseteq M$ é aberto em $M$.

    Pela proposição anterior,
    \[
        \lim_{x\to p} f(x) = \lim_{x\to p} f|_A(x) = f(p),
    \]

    Portanto, $f$ é contínua em $p$.
\end{proof}

Por fim, a noção de limite não depende do que ocorre no ponto para o qual o limite é calculado.


\begin{proposition}Sejam
    \begin{itemize}
        \item $f: M \to N$ uma função,
        \item $p\in M$ é ponto de acumulação de $M$,
        \item $f: M \to N$ uma função, e
        \item $L\in N$.
    \end{itemize}

    Então $p$ é um ponto de acumulação de $M$ e $L$ é limite de $f$ em $p$, se, e somente se $L$ é limite de $f|_{M\setminus\{p\}}$ em $p$.

    Em outras palavras, nesse caso,
    \[
        \lim_{x\to p} f|_{M\setminus\{p\}}(x) = \lim_{x\to p} f(x).
    \]
\end{proposition}

\begin{proof}
    Já vimos que $M$ e $M\setminus\{p\}$ têm os mesmos pontos de acumulação, pois $\{p\}$ é finito.

    Se $L$ é limite de $f$ em $p$, então, $L$ é limite de $f|_{M\setminus\{p\}}$ em $p$, pois limites são preservados por restrições.

    Reciprocamente, se $L$ é limite de $f|_{M\setminus\{p\}}$ em $p$, temos que
    \[
        \forall \epsilon > 0\, \exists \delta > 0\, \forall x \in M\setminus\{p\}, (d(x, p) < \delta \implies d'(f|_{M\setminus\{p\}}(x), L) < \epsilon).
    \]
    O que é imediatamente equivalente a
    \[
        \forall \epsilon > 0\, \exists \delta > 0\, \forall x \in M\setminus\{p\}, (d(x, p) < \delta \implies d'(f(x), L) < \epsilon).
    \]

    Assim, $L$ é limite de $f$ em $p$.
\end{proof}


\begin{example}\label{ex:restricaoContinuidade}
    Seja $f:\mathbb R^2\rightarrow \mathbb R$ dada por
    \[
        f(x, y) = \begin{cases}
            1, & \text{se } (x, y) \neq (0, 0),\\
            0, & \text{se } (x, y) = (0, 0).
        \end{cases}
    \]

    Temos que $f$ é contínua nos pontos de $A=\mathbb R^2\setminus\{(0, 0)\}$, pois $f|_{A}$ é uma função constante, que é contínua, e $A$ é um conjunto aberto.

    No entanto, $f$ não é contínua em $(0, 0)$.
    De fato, primeiro note que $(0, 0)$.
    Como o valor de $f$ em $(0, 0)$ não importa para o limite, temos
    \[
        \lim_{(x, y)\to (0, 0)} f(x, y) = \lim_{(x, y)\to (0, 0)} f|_{A}(x, y) = \lim_{(x, y)\to (0, 0)} 1|_{A}(x, y) = \lim_{(x, y)\to (0, 0)} 1|_{\mathbb R^2} = 1,
    \]
    em que $1_A$ e $1_{\mathbb R^2}$ são as funções constantes de $A$ e $\mathbb R^2$ identicamente iguais a $1$, respectivamente.

    Assim, $\lim_{(x, y)\to (0, 0)} f(x, y) = 1 \neq f(0, 0)$, e, portanto, $f$ não é contínua em $(0, 0)$.
\end{example}

Uma outra questão que pode surgir naturalmente é a seguinte: sabemos que a função $f:\mathbb R\rightarrow \mathbb R$ dada por $f(x) = x^2$ é contínua.
Porém, essa mesma função pode ser vista como uma função de $\mathbb R$ em $[0, +\infty)$, e, pode-se perguntar, se nesse sentido, ela ainda é contínua.

A proposição abaixo responde essa pergunta positivamente, e de forma geral.

\begin{proposition}Sejam
    \begin{itemize}
        \item $f: M \to N$ uma função,
        \item $a \in M$ um ponto de acumulação de $M$.
        \item $N'$ um subespaço de $N$ tal que $f[M]\subseteq N'$,
        \item $L\in N'$
    \end{itemize}
    Então $L$ é limite de $f:M\to N$ em $a$, se, e somente se $L$ é limite de $f:M\to N'$ em $a$.
    Em outras palavras, restringir contradomínios não altera a noção de limite para pontos do espaço restrito.
\end{proposition}
\begin{proof}
    Suponha que $L$ é limite de $f:M\to N'$ em $a$.
    Seja $i: N'\rightarrow N$ a inclusão de $N'$ em $N$.
    Então $L=i(L)\in M$ é limite de $i\circ f:M\to N$ em $a$, pois $i$ é contínua.
    Porém, $i\circ f=f$, assim, segue a tese.

    Reciprocamente, suponha que $L$ é limite de $f:M\to N$ em $a$.
    Para ver que $L$ é limite de $f:M\to N'$ em $a$, devemos ver que dado $\epsilon>0$, existe $\delta>0$ tal que, para todo $x\in M\setminus\{a\}$, se $d(x, a) < \delta$, então $d'(f(x), L) < \epsilon$.
    Mas isso é diretamente garantido pelo fato de $L$ ser limite de $f:M\to N$ em $a$.
\end{proof}

\begin{example}
    Seja $f:\mathbb R^2\rightarrow \mathbb R^2\setminus\{(0, 0)\}$ dada por
    \[
        f(x, y) =\begin{cases}
            (x, y), & \text{se } (x, y) \neq (0, 0),\\
            (1, 1), & \text{se } (x, y) = (0, 0).
        \end{cases}
    \]

    Então $f:\mathbb R^2\rightarrow \mathbb R^2\setminus\{(0, 0)\}$ não possui limite considerando o contradomínio $\mathbb R^2\setminus\{(0, 0)\}$.
    
    De fato, se possuísse, existiria $L\in \mathbb R^2\setminus\{(0, 0)\}$ tal que $L$ é limite de $f:\mathbb R^2\rightarrow \mathbb R^2\setminus\{(0, 0)\}$.
    Porém, $L$ seria limite de $f:\mathbb R^2\rightarrow \mathbb R^2$ em $(0, 0)$, o que é impossível, pois teríamos:
    \begin{align*}
            & L \text{ é limite de } f \text{ em } (0, 0) \\
        \implies & L \text{ é limite de } f|_{\mathbb R^2\setminus\{(0, 0)\}} \text{ em } (0, 0)\\
        \implies & L \text{ é limite de } \id_{\mathbb R^2}|_{\mathbb R^2\setminus\{(0, 0)\}} \text{ em } (0, 0)\\
        \implies & L \text{ é limite de } \id_{\mathbb R^2} \text{ em } (0, 0)\\
        \implies & L = (0, 0),
    \end{align*}
    E, portanto, $(0, 0)=L\in \mathbb R^2\setminus\{(0, 0)\}$.
\end{example}

\begin{corollary}
    Sejam
    \begin{itemize}
        \item $f: M \to N$ uma função,
        \item $a \in M$.
        \item $N'$ um subespaço de $N$ tal que $f[M]\subseteq N'$,
    \end{itemize}
    Então $f:M\to N$ é contínua em $a$, se, e somente se, $f:M\to N'$ é contínua em $a$.
\end{corollary}
\begin{proof}
    Se $a$ é um ponto isolado de $M$, então ambas $f:M\to N$ e $f:M\to N'$ são contínuas em $a$ trivialmente.
    Se $a$ é um ponto de acumulação de $M$, então:
    \begin{align*}
            & f:M\to N \text{ é contínua em } a\\
        \iff & f(a) \text{ é limite de } f:M\to N \text{ em } a\\
        \iff & f(a) \text{ é limite de } f:M\to N' \text{ em } a\\
        \iff & f:M\to N' \text{ é contínua em } a.
    \end{align*}
\end{proof}